\chapter{Metodologia}
\label{ch:metodologia}

A presente investigação adota uma abordagem
\textbf{quantitativa e experimental}, fundamentada na modelagem
matemática de sistemas de software através da Teoria dos Grafos e na
execução repetível do \emph{pipeline} analítico sobre projetos reais
de código aberto. O Capítulo~\ref{ch:metodologia} detalha:
Secção~\ref{sec:estrategia} --- a estratégia de investigação;
Secção~\ref{sec:reqgraph} --- a arquitetura do protótipo \emph{ReqGraph};
Secção~\ref{sec:tecnologia} --- as escolhas de implementação;
Secção~\ref{sec:protocolo} --- o corpus experimental, o protocolo
\code{func\_to\_req} e os critérios de avaliação. A Secção~\ref{sec:metricas}
formaliza as métricas de comparação de rankings.

\noindent\textbf{Conceitos-chave:} abordagem quantitativa; análise
estática; AST; \emph{pipeline} reprodutível; protocolo experimental;
\emph{rank correlation}.

% ---------------------------------------------------------------------------
\section{Estratégia de Investigação}
\label{sec:estrategia}

A investigação está estruturada em quatro etapas. Em primeiro lugar,
\textbf{fundamentou-se teoricamente} a aplicação de algoritmos de
ranqueamento de grafos à Priorização de Requisitos de Software (PRS),
conforme exposto no Capítulo~\ref{ch:fundamentos}, consolidando o
entendimento sobre as diferenças estruturais entre o HITS
(dualidade \emph{hubs}/autoridades), o SALSA (passeios aleatórios em
grafos bipartidos) e o PageRank (distribuição estacionária em cadeias
de Markov). Em segundo lugar, \textbf{projetou-se e implementou-se}
um protótipo --- o \emph{ReqGraph} --- capaz de extrair
automaticamente o grafo de dependências entre requisitos a partir do
código-fonte de uma aplicação Python, recorrendo a análise estática
via \emph{Abstract Syntax Tree} (AST). Em terceiro lugar,
\textbf{definiu-se um protocolo experimental} para aplicar os
algoritmos PageRank, HITS e SALSA aos grafos extraídos de três
projetos de código aberto de complexidade crescente. Em quarto lugar,
\textbf{quantificou-se a concordância} entre rankings via Spearman
$\rho$ e Kendall $\tau_b$, respondendo a RQ1, RQ2 e RQ3
(Secção~\ref{sec:objetivos}).

O ambiente de experimentação é \textbf{local}, em Python 3.8+,
recorrendo às bibliotecas \code{ast} (biblioteca padrão),
\code{NetworkX} (manipulação de grafos), \code{NumPy}/\code{SciPy}
(computação matricial) e \code{Matplotlib} (visualização). Esta
escolha justifica-se pela necessidade de manipular matrizes de
adjacência e realizar cálculos iterativos de autovetores de forma
escalável, em ambiente totalmente reprodutível e sem dependência de
serviços externos.

A opção por \textbf{análise estática} (em detrimento de instrumentação
dinâmica) deve-se à reprodutibilidade dos resultados: o grafo
extraído depende exclusivamente do código-fonte, sem variabilidade
introduzida por entradas de execução. Esta opção é coerente com a
literatura sobre rastreabilidade de requisitos, que defende a
derivação determinística do grafo de dependências a partir de
artefactos estáveis do sistema
\parencite{silva2018, gotel1994, rasiman2022}.

A análise dos resultados focar-se-á em três dimensões:
(i) \textbf{Consistência Algorítmica} --- comparação entre os
rankings gerados pelos três algoritmos, recorrendo a coeficientes de
correlação de Spearman e Kendall (RQ1);
(ii) \textbf{Caracterização Estrutural} --- densidade, presença de
ciclos e identificação dos vértices centrais em cada Grafo de
Requisitos (RQ2);
(iii) \textbf{Coerência Arquitetural} --- comparação qualitativa
entre os requisitos ranqueados como mais importantes e a arquitetura
conhecida de cada projeto (RQ3).

% ---------------------------------------------------------------------------
\section{Arquitetura do Protótipo ReqGraph}
\label{sec:reqgraph}

O \emph{ReqGraph} é o componente que materializa a transformação do
código-fonte num \textbf{Grafo de Requisitos}
$G_R = (V_R, E_R)$, em que cada vértice $v \in V_R$ representa um
requisito de domínio e cada aresta direcionada
$(r_i, r_j) \in E_R$ traduz uma dependência de implementação
detetada entre os requisitos $r_i$ e $r_j$. O \emph{pipeline}
analítico é composto por quatro estágios sequenciais:

\begin{enumerate}
    \item \textbf{Análise estática via AST.} Cada ficheiro Python
        (\code{.py}) do projeto é parseado pelo módulo \code{ast} da
        biblioteca padrão. O resultado é uma árvore sintática a
        partir da qual são extraídas definições de funções
        (\code{FunctionDef}), métodos de classes e invocações
        (\code{Call}).
    \item \textbf{Construção do Call Graph.} As invocações detetadas
        dão origem a um grafo direcionado $G_C = (V_C, E_C)$ em que
        cada vértice corresponde a uma função (ou método) e cada
        aresta $(f_i, f_j)$ indica que $f_i$ invoca $f_j$. O
        algoritmo resolve referências entre módulos (via
        \code{import}), referências a métodos da própria classe
        (\code{self.method()}) e qualifica funções pelo seu
        \emph{namespace} para evitar colisões de nomes.
    \item \textbf{Mapeamento \code{func\_to\_req}.} É fornecido pelo
        investigador um dicionário que associa cada função a um
        requisito de domínio (\eg\
        \code{parse\_args $\to$ REQ\_PARSING}). Este artefacto é o
        ponto de contacto entre o nível técnico (funções) e o nível
        semântico (requisitos). Para reduzir a fricção da sua
        produção, foi também desenvolvido um \emph{prompt}
        estruturado (\code{reqgraph/llm\_prompt\_mapping.md}) que
        permite gerar o mapeamento de forma assistida via modelos de
        linguagem de grande escala (ChatGPT, Claude, Gemini). O
        protocolo de construção deste mapeamento é descrito em
        detalhe na Secção~\ref{sec:protocolo-func2req}.
    \item \textbf{Derivação do Grafo de Requisitos.} Para cada
        aresta $(f_i, f_j) \in E_C$ tal que $f_i$ está mapeada para
        o requisito $r_a$ e $f_j$ está mapeada para o requisito
        $r_b$, com $r_a \neq r_b$, é introduzida a aresta
        $(r_a, r_b)$ em $E_R$. Arestas duplicadas e auto-laços
        (correspondentes a invocações internas dentro do mesmo
        requisito) são descartados.
\end{enumerate}

\begin{figure}[ht]
    \centering
    \includegraphics[width=0.92\linewidth]{call_graph_simples_stdlib_excerto.png}
    \caption{Excerto representativo do \emph{Call Graph} extraído via
    AST do CPython \emph{stdlib} (\code{argparse}). Cada vértice é uma
    função ou método --- e o seu nome está rotulado no vértice ---,
    enquanto cada aresta $(f_i, f_j)$ indica que $f_i$ invoca $f_j$. As
    setas escuras terminam na borda do nó, em coerência com a
    convenção visual dos grafos do Capítulo~\ref{ch:fundamentos}.}
    \label{fig:call-graph-stdlib}
\end{figure}

Por legibilidade, a Figura~\ref{fig:call-graph-stdlib} apresenta apenas
um \emph{excerto} do \emph{Call Graph} do caso \emph{stdlib} (o grafo
completo tem $120$ vértices e $130$ arestas). De modo a documentar a
escala real de cada caso de estudo sem sacrificar a legibilidade dos
nomes, o \emph{Call Graph} dos três projetos é apresentado, no
Apêndice~\ref{ap:callgraph-repr}, em três representações
complementares: (i) o grafo completo rotulado; (ii) a sua maior
componente conexa; e (iii) um excerto representativo.

O código-fonte completo do \emph{ReqGraph}, incluindo o módulo
\code{ranker.py} descrito na Secção~\ref{sec:tecnologia}, está
disponível em \url{https://github.com/tales96323/Tales_Tese_ulusofona}.

% ---------------------------------------------------------------------------
\section{Implementação e Tecnologias}
\label{sec:tecnologia}

O protótipo foi implementado em \textbf{Python 3.8+}, organizado como
pacote instalável via \code{pip install -e reqgraph/}. As escolhas
tecnológicas seguem critérios de robustez e adesão a \emph{standards}
da comunidade científica (Tabela~\ref{tab:tecnologias}).

\begin{table}[ht]
    \centering
    \caption{Tecnologias utilizadas e respetivas funções no
    \emph{pipeline}.}
    \label{tab:tecnologias}
    \begin{tabular}{@{}lll@{}}
        \toprule
        \textbf{Tecnologia} & \textbf{Versão mín.} & \textbf{Função no \emph{pipeline}} \\
        \midrule
        \code{ast} (biblioteca padrão)
            & ---            & \emph{Parsing} sintático do código-fonte Python \\
        \code{networkx} & $\geq 3.0$ & Manipulação e análise de grafos direcionados \\
        \code{matplotlib} & $\geq 3.5$ & Geração de visualizações em PNG \\
        \code{numpy} / \code{scipy} & $\geq 1.24$ / $\geq 1.10$ & Computação matricial \\
        \code{graphviz} (DOT) & --- & Formato de exportação canónico \\
        \bottomrule
    \end{tabular}
\end{table}

A interface de linha de comando
(\code{python -m reqgraph <projeto> --mapping <mapeamento.py>}) expõe
o \emph{pipeline} completo, gerando como artefactos:
\code{call\_graph.png}, \code{req\_graph.png},
\code{req\_graph.json} (lista de adjacência) e
\code{req\_graph.dot} (formato Graphviz).

\subsection{Módulo de Ranqueamento (\code{ranker.py})}
\label{sec:ranker}

O módulo \code{ranker.py} consome o ficheiro \code{req\_graph.json}
produzido pelo \emph{ReqGraph} e aplica três algoritmos de
ranqueamento sobre o grafo direcionado $G_R$. A semântica adotada
para as arestas é: $r_a \rightarrow r_b$ significa que o requisito
$r_a$ \textbf{depende de} $r_b$. Em consequência, o \emph{in-degree}
de um vértice mede quantos requisitos dele dependem (uma autoridade
no sentido de Kleinberg), enquanto o \emph{out-degree} mede quantos
outros requisitos ele invoca (um \emph{hub}).

\subsection{PageRank}
\label{sec:impl-pagerank}

A implementação recorre a \code{nx.pagerank} com parâmetros
canónicos: fator de amortecimento $d = 0{,}85$, número máximo de
iterações $100$ e tolerância de convergência $10^{-6}$.

\paragraph{Tratamento de \emph{Dangling Nodes}.}
\label{sec:dangling}
Vértices sem arestas de saída são tratados internamente pelo
NetworkX por \textbf{redistribuição uniforme}: a massa de
probabilidade que um vértice sem arestas de saída teria de
``exportar'' é dividida em partes iguais entre todos os vértices do
grafo, mantendo a matriz de transição estocástica. Este artifício é
necessário para garantir que a cadeia de Markov subjacente é
ergódica e, portanto, que a iteração converge para uma única
distribuição estacionária \parencite{brin1998}. A soma dos
\emph{scores} totaliza $1{,}0$, permitindo interpretá-los como uma
distribuição de probabilidade.

\subsection{HITS}
\label{sec:impl-hits}

A implementação recorre a \code{nx.hits} com tolerância $10^{-8}$ e
$100$ iterações máximas. Retorna dois vetores normalizados (norma
euclidiana unitária): \emph{hubs} e \emph{authorities}. A
interpretação para a PRS é direta: um requisito com elevado
\emph{authority} é uma \textbf{dependência crítica} do sistema
(alterações nele propagam-se a muitos consumidores), enquanto um
requisito com elevado \emph{hub} é um \textbf{orquestrador}
(alterações nele afetam muitas integrações).

\subsection{SALSA}
\label{sec:impl-salsa}

O NetworkX \textbf{não fornece} uma implementação do SALSA, pelo que
esta foi desenvolvida manualmente seguindo \citet{lempel2001}. Para
cada aresta $(u, v) \in E_R$, constroem-se duas matrizes de
transição:

\begin{itemize}
    \item $H \rightarrow A$: probabilidade $1/\text{out}(u)$ de um
        \emph{hub} $u$ transitar para a \emph{authority} $v$.
    \item $A \rightarrow H$: probabilidade $1/\text{in}(v)$ de uma
        \emph{authority} $v$ regressar a um \emph{hub} $u$.
\end{itemize}

A cadeia de \emph{hubs} tem matriz
$T_{\text{hub}} = (A \rightarrow H)(H \rightarrow A)$, e a de
\emph{authorities} tem
$T_{\text{auth}} = (H \rightarrow A)(A \rightarrow H)$. Os
\emph{scores} finais são as distribuições estacionárias obtidas por
\emph{power iteration} com tolerância $10^{-8}$. Linhas nulas
(correspondentes a vértices absorventes) são tratadas por
redistribuição uniforme, em analogia ao tratamento de
\emph{dangling nodes} do PageRank.

A diferença substantiva face ao HITS é a \textbf{normalização local}
(por grau de cada vértice), que torna o SALSA teoricamente menos
sensível ao \textbf{Efeito TKC}. Como discutido em
Secção~\ref{sec:salsa}, esta vantagem só é observável empiricamente
em grafos com sub-comunidades coesas embutidas; o corpus utilizado
nesta dissertação não permite, por dimensão, observar este fenómeno.

\subsection{Artefactos Gerados}
\label{sec:artefactos}

Para cada projeto analisado, o \code{ranker.py} produz dois
artefactos: \code{ranking\_results.json} (com os \emph{scores}
completos dos três algoritmos) e \code{ranking\_results.png} (com
gráficos de barras horizontais agrupados por algoritmo). Quando
executado com a \emph{flag} \code{--all}, gera adicionalmente
\code{ranking\_consolidado.json} e \code{ranking\_comparativo.png}
na raiz do projeto, agregando os resultados dos três casos de estudo.
A partir de R5, é também produzido \code{rank\_correlations.json}
pelo \emph{script} auxiliar \code{compute\_rank\_correlations.py},
que contém os coeficientes Spearman e Kendall para cada par de
algoritmos em cada caso de estudo.

% ---------------------------------------------------------------------------
\section{Métricas de Comparação de Rankings}
\label{sec:metricas}

A comparação entre rankings produzidos por algoritmos diferentes ---
elemento central de RQ1 e RQ2 --- é feita de forma formal,
recorrendo a duas métricas clássicas de \emph{rank correlation}:

\begin{description}
    \item[Coeficiente de Spearman ($\rho$).] Pearson aplicado aos
        \emph{ranks} (em vez dos valores brutos), com tratamento de
        \emph{ties} via \emph{mean ranks}. Mede correlação monotónica
        \parencite{spearman1904}.
        \[
            \rho = 1 - \frac{6 \sum_i d_i^2}{n(n^2 - 1)}
            \qquad \text{(sem \emph{ties})}
        \]
        Com \emph{ties}, é calculado como a correlação de Pearson dos
        \emph{mean ranks} (definição usada no script
        \code{compute\_rank\_correlations.py}).

    \item[Coeficiente de Kendall ($\tau_b$).] Baseado em pares
        concordantes e discordantes \parencite{kendall1938};
        $\tau_b$ é a variante com correção de \emph{ties}
        \parencite{goodman1954}:
        \[
            \tau_b = \frac{|C| - |D|}
                {\sqrt{(P_0 - T_a)(P_0 - T_b)}}
        \]
        onde $|C|$ e $|D|$ são os números de pares concordantes e
        discordantes, $P_0 = n(n-1)/2$ e $T_a, T_b$ contam pares com
        \emph{ties} em cada ranking.
\end{description}

\paragraph{Critério a priori de ``convergência''.} Para fins de
discussão (Secção~\ref{sec:discussao}), define-se \emph{a priori},
antes de observar os dados, que existe \textbf{convergência entre dois
algoritmos} quando se verifica qualquer um dos seguintes critérios:

\begin{enumerate}[label=(\alph*)]
    \item Spearman $\rho \geq 0{,}7$ \emph{e} Kendall $\tau_b \geq 0{,}55$
        (ambos os limiares correspondem ao limiar usual de
        ``\emph{strong correlation}'' na literatura \parencite{wohlin2012});
    \item \emph{Top-1} idêntico em pelo menos dois dos três
        algoritmos.
\end{enumerate}

Quando nenhum dos critérios se verifica, fala-se de
\textbf{divergência}; nesse caso, a interpretação é qualitativa (cada
algoritmo capta uma dimensão diferente de centralidade, conforme
discussão em Secção~\ref{sec:discussao}).

% ---------------------------------------------------------------------------
\section{Corpus Experimental e Protocolo de Avaliação}
\label{sec:protocolo}

\subsection{Seleção dos Casos de Estudo}
\label{sec:selecao}

Para validar a abordagem em diferentes regimes de complexidade
arquitetural, foram selecionados \textbf{três projetos reais de
código aberto} disponíveis no GitHub, organizados por nível
crescente de acoplamento esperado
(Tabela~\ref{tab:projetos}). A escolha procura cobrir
\emph{regimes distintos}: biblioteca padrão monolítica
(CPython \emph{stdlib}), \emph{microframework} web (Flask) e
\emph{framework} de \emph{machine learning} com forte composição
funcional (\emph{scikit-learn}). A seleção privilegia ficheiros
representativos do \emph{core} arquitetural de cada projeto,
evitando módulos auxiliares (testes, utilitários, \emph{type stubs}).

\begin{table}[ht]
    \centering
    \caption{Projetos selecionados como casos de estudo.}
    \label{tab:projetos}
    \footnotesize
    \begin{tabularx}{\linewidth}{@{}l l X l@{}}
        \toprule
        \textbf{Nível} & \textbf{Projeto} & \textbf{Ficheiros Analisados} & \textbf{Origem (GitHub)} \\
        \midrule
        Simples & CPython \emph{stdlib}
            & \code{Lib/argparse.py}, \code{Lib/http/server.py}
            & \url{https://github.com/python/cpython} \\
        Médio & Flask
            & \code{src/flask/app.py}, \code{src/flask/cli.py},
              \code{src/flask/blueprints.py}
            & \url{https://github.com/pallets/flask} \\
        Complexo & \emph{scikit-learn}
            & \code{sklearn/pipeline.py},
              \code{sklearn/linear\_model/\_base.py},
              \code{sklearn/linear\_model/\_logistic.py},
              \code{sklearn/tree/\_classes.py}
            & \url{https://github.com/scikit-learn/scikit-learn} \\
        \bottomrule
    \end{tabularx}
\end{table}

\subsection{Protocolo de Construção do Mapeamento
\code{func\_to\_req}}
\label{sec:protocolo-func2req}

Dado que o passo (3) do \emph{pipeline}
(Secção~\ref{sec:reqgraph}) é a principal fonte de subjetividade
residual, o presente trabalho formaliza explicitamente um protocolo
para a sua construção:

\begin{enumerate}
    \item \textbf{Levantamento de domínios candidatos.} Identificar,
        a partir da documentação oficial do projeto, módulos
        principais e \emph{API} pública, um conjunto inicial de
        \emph{features} de alto nível (\eg\ \emph{routing},
        \emph{request handling}, \emph{error handling} no Flask).
    \item \textbf{Definição de granularidade.} Cada requisito deve
        agregar entre 5 e 30 funções; conjuntos com menos de 5
        funções fundem-se com domínios próximos; conjuntos com mais
        de 30 dividem-se segundo a sub-estrutura modular.
    \item \textbf{Atribuição.} Cada função do código-fonte é
        atribuída a exatamente um requisito; funções utilitárias
        partilhadas por vários requisitos são atribuídas ao
        requisito do qual são chamadas mais frequentemente.
    \item \textbf{Revisão de cobertura.} Verificar que todas as
        funções identificadas pela AST estão mapeadas; funções não
        mapeadas tornam-se \emph{dangling nodes} no \emph{Call Graph}
        e influenciam o ranking via redistribuição uniforme.
    \item \textbf{Validação de coerência.} Inspecionar
        manualmente o \code{req\_graph.png} gerado: arestas
        inesperadas (\eg\ ciclos que contradigam a arquitetura
        conhecida) podem indicar erros de mapeamento; arestas
        ausentes (\eg\ ausência de dependências esperadas) podem
        indicar requisitos demasiado granulares.
\end{enumerate}

\begin{notebox}[Limitação reconhecida]
    A presente dissertação executa este protocolo com um único
    investigador, sem \emph{inter-rater reliability} (\eg\
    Cohen's $\kappa$). A replicação por um segundo investigador,
    com cálculo formal da concordância, fica identificada como
    trabalho futuro (Secção~\ref{sec:trabalho-futuro}) e como
    \emph{ameaça de validade interna}
    (Secção~\ref{sec:ameacas}).
\end{notebox}

\subsection{Protocolo Experimental}
\label{sec:protocolo-exp}

Cada caso de estudo foi processado seguindo os mesmos seis passos,
automatizados pelo \emph{script} \code{run\_tests.py}:

\begin{enumerate}
    \item Cópia dos ficheiros-alvo do repositório de origem para a
        pasta \code{testes/<nivel>/}.
    \item Construção manual do dicionário \code{func\_to\_req} em
        \code{mapeamento.py}, seguindo o protocolo da
        Secção~\ref{sec:protocolo-func2req}.
    \item Execução do \emph{ReqGraph}:\\
        \code{python -m reqgraph testes/<nivel>/ --mapping testes/<nivel>/mapeamento.py}
    \item Inspeção dos artefactos gerados
        (\code{call\_graph.png}, \code{req\_graph.png},
        \code{req\_graph.json}).
    \item Execução do \code{ranker.py} sobre o
        \code{req\_graph.json} gerado.
    \item Cálculo dos coeficientes de Spearman e Kendall entre cada
        par de algoritmos, via \code{compute\_rank\_correlations.py}.
\end{enumerate}

\subsection{Dimensões de Análise}
\label{sec:dimensoes}

A análise dos resultados (Capítulo~\ref{ch:resultados}) é
estruturada em três dimensões:

\begin{itemize}
    \item \textbf{Resultados Técnicos} --- métricas estruturais do
        \emph{pipeline}: número de funções analisadas, arestas no
        \emph{Call Graph}, requisitos identificados, arestas no
        Grafo de Requisitos e densidade.
    \item \textbf{Resultados Analíticos} --- interpretação dos
        rankings: identificação dos requisitos mais críticos,
        comparação entre os três algoritmos com Spearman e Kendall,
        e discussão à luz da arquitetura conhecida de cada projeto.
    \item \textbf{Ameaças à Validade} --- categorização das
        limitações segundo o referencial de
        \citet{wohlin2012}: \emph{construct}, \emph{internal},
        \emph{external} e \emph{conclusion}.
\end{itemize}
